% !TeX program = xelatex
% Overleaf: Menu -> Compiler -> XeLaTeX
\documentclass[11pt,a4paper,oneside]{report}

\usepackage[a4paper,top=2.25cm,bottom=2.25cm,left=2.35cm,right=2.35cm,headheight=15pt]{geometry}
\usepackage{fontspec}
\usepackage{polyglossia}
\setdefaultlanguage{turkish}
\setotherlanguage{english}
\defaultfontfeatures{Ligatures=TeX,Scale=MatchLowercase}
\setmainfont{TeX Gyre Pagella}
\setsansfont{TeX Gyre Heros}
\setmonofont{DejaVu Sans Mono}[Scale=0.86]

\usepackage{microtype}
\usepackage{setspace}
\usepackage{parskip}
\usepackage{ragged2e}
\usepackage{enumitem}
\usepackage{amsmath,amssymb}
\usepackage{booktabs}
\usepackage{longtable}
\usepackage{tabularx}
\usepackage{array}
\usepackage{makecell}
\usepackage{graphicx}
\usepackage{caption}
\usepackage{float}
\usepackage[most]{tcolorbox}
\usepackage{xcolor}
\usepackage{fancyhdr}
\usepackage{titlesec}
\usepackage{hyperref}
\usepackage{bookmark}
\usepackage{xurl}

\definecolor{Navy}{HTML}{16324F}
\definecolor{Blue}{HTML}{245B8A}
\definecolor{LightBlue}{HTML}{EAF3FA}
\definecolor{DarkRed}{HTML}{8B1E1E}
\definecolor{LightRed}{HTML}{FBECEC}
\definecolor{DarkGreen}{HTML}{276749}
\definecolor{LightGreen}{HTML}{EAF7F0}
\definecolor{DarkGold}{HTML}{806000}
\definecolor{LightGold}{HTML}{FFF8DB}
\definecolor{MidGray}{HTML}{5F6B76}
\definecolor{LightGray}{HTML}{F4F6F8}
\definecolor{RuleGray}{HTML}{D6DCE2}

\hypersetup{
  colorlinks=true,
  linkcolor=Navy,
  urlcolor=Blue,
  citecolor=DarkGreen,
  pdftitle={İHA ve ADS-B Telemetrisi Anomali Tespiti: Kapsamlı Teknik Rapor ve Fizibilite Hükmü},
  pdfauthor={AnomalyDetection Projesi},
  pdfsubject={Makine öğrenmesi, fiziksel tutarlılık, derin öğrenme ve operasyonel fizibilite değerlendirmesi},
  pdfkeywords={anomaly detection, ADS-B, UAV, machine learning, LSTM, Isolation Forest, CUSUM, feasibility}
}

\setstretch{1.08}
\setlist[itemize]{leftmargin=1.55em,itemsep=0.22em,topsep=0.35em}
\setlist[enumerate]{leftmargin=1.75em,itemsep=0.28em,topsep=0.35em}
\setlength{\emergencystretch}{3em}
\setlength{\LTpre}{0.5em}
\setlength{\LTpost}{0.8em}
\renewcommand{\arraystretch}{1.16}

\newcolumntype{Y}{>{\RaggedRight\arraybackslash}X}
\newcolumntype{L}[1]{>{\RaggedRight\arraybackslash}p{#1}}
\newcommand{\code}[1]{\texttt{\detokenize{#1}}}
\newcommand{\metric}[1]{\textbf{\textcolor{Navy}{#1}}}
\newcommand{\pass}{\textbf{\textcolor{DarkGreen}{GEÇTİ}}}
\newcommand{\fail}{\textbf{\textcolor{DarkRed}{KALDI/BAŞARISIZ}}}
\newcommand{\na}{\textcolor{MidGray}{uygulanmadı}}
\newcommand{\fasaat}{yanlış alarm/saat}
\newcommand{\FA}{FA}

\newtcolorbox{verdictbox}{
  colback=LightRed,
  colframe=DarkRed,
  boxrule=0.9pt,
  arc=2mm,
  left=4mm,right=4mm,top=3mm,bottom=3mm,
  breakable
}

\newtcolorbox{infobox}[1][]{
  colback=LightBlue,
  colframe=Blue,
  boxrule=0.6pt,
  arc=2mm,
  left=4mm,right=4mm,top=2.5mm,bottom=2.5mm,
  title=#1,
  fonttitle=\bfseries\sffamily,
  breakable
}

\newtcolorbox{warningbox}[1][]{
  colback=LightGold,
  colframe=DarkGold,
  boxrule=0.6pt,
  arc=2mm,
  left=4mm,right=4mm,top=2.5mm,bottom=2.5mm,
  title=#1,
  fonttitle=\bfseries\sffamily,
  breakable
}

\newtcolorbox{successbox}[1][]{
  colback=LightGreen,
  colframe=DarkGreen,
  boxrule=0.6pt,
  arc=2mm,
  left=4mm,right=4mm,top=2.5mm,bottom=2.5mm,
  title=#1,
  fonttitle=\bfseries\sffamily,
  breakable
}

\titleformat{\chapter}[display]
  {\normalfont\sffamily\bfseries\color{Navy}}
  {\filleft\Large\chaptername\ \thechapter}
  {0.5ex}
  {\titlerule\vspace{1ex}\Huge}
  [\vspace{0.5ex}\titlerule]

\titleformat{\section}
  {\Large\sffamily\bfseries\color{Navy}}
  {\thesection}{0.7em}{}

\titleformat{\subsection}
  {\large\sffamily\bfseries\color{Blue}}
  {\thesubsection}{0.7em}{}

\titleformat{\subsubsection}
  {\normalsize\sffamily\bfseries\color{MidGray}}
  {\thesubsubsection}{0.7em}{}

\pagestyle{fancy}
\fancyhf{}
\fancyhead[L]{\small\sffamily İHA/ADS-B Anomali Tespiti}
\fancyhead[R]{\small\sffamily Teknik Fizibilite Raporu}
\fancyfoot[C]{\small\sffamily \thepage}
\renewcommand{\headrulewidth}{0.4pt}
\renewcommand{\footrulewidth}{0pt}

\begin{document}

\begin{titlepage}
  \centering
  \vspace*{1.3cm}

  {\sffamily\bfseries\color{Navy}\fontsize{25}{31}\selectfont
  İHA ve ADS-B Telemetrisi\\[0.25em]
  Anomali Tespiti\par}

  \vspace{0.8cm}

  {\sffamily\color{Blue}\fontsize{16}{21}\selectfont
  Kapsamlı Makine Öğrenmesi Raporu\\
  ve Teknik Fizibilite Hükmü\par}

  \vspace{1.4cm}

  \begin{tcolorbox}[
    width=0.88\textwidth,
    colback=LightGray,
    colframe=RuleGray,
    boxrule=0.6pt,
    arc=2mm,
    left=5mm,right=5mm,top=4mm,bottom=4mm
  ]
  \centering
  \sffamily
  \textbf{Kapsam}\\[0.3em]
  Arşivlenmiş non-ADS-B çalışmaları (ML-0--ML-16), RFLY çalışmaları,
  reddedilmiş ilk ADS-B yaklaşımları ve aktif ADS-B contextual-physics hattı
  birlikte değerlendirilmiştir.
  \end{tcolorbox}

  \vfill

  \begin{tabular}{rl}
    \textbf{Rapor tarihi:} & 16 Temmuz 2026\\
    \textbf{Proje deposu:} & \code{AnomalyDetection}\\
    \textbf{Belge niteliği:} & Teknik kapanış ve yönetim değerlendirmesi\\
    \textbf{Derleyici:} & XeLaTeX\\
  \end{tabular}

  \vspace{1.2cm}

  {\small\color{MidGray}
  Bu rapor, deneylerin olumlu ve olumsuz sonuçlarını aynı ölçüm disipliniyle
  sunmak; ``neyi, neden kullandık; sorunu neden çözmedi; hangi düzeltmeleri
  yaptık; neden yine yeterli olmadı?'' sorularına kanıtlı cevap vermek amacıyla
  hazırlanmıştır.\par}
\end{titlepage}

\pagenumbering{roman}

\chapter*{Belge Durumu ve Ana Hüküm}
\addcontentsline{toc}{chapter}{Belge Durumu ve Ana Hüküm}

\begin{verdictbox}
\textbf{\large Teknik hüküm}

Mevcut veri setleri, etiket sözleşmeleri, sensör gözlenebilirliği, platformlar
arası dağılım farkı ve kabul edilen operasyonel yanlış alarm bütçesi altında;
farklı platform ve anomali ailelerini kapsayan, düşük yanlış alarmlı,
operasyonel kullanıma hazır \textbf{genel bir anomali tespit sistemi
achievable değildir}. Projenin mevcut kapsamıyla ürün fizibilitesi
doğrulanmamış; deneysel kanıtlar hedefin bu koşullarda karşılanamadığını
göstermiştir.

\medskip

Bu hüküm, ``anomali tespiti hiçbir koşulda mümkün değildir'' anlamına gelmez.
Dar kapsamlı fiziksel limit denetleyicileri, veri-kalitesi monitörleri,
belirli sensör veya arıza türlerine özel dedektörler ve araştırma prototipleri
yapılabilir. Ancak bunlar başlangıçta hedeflenen genel ve operasyonel ML
ürünüyle aynı iddiayı taşımaz.
\end{verdictbox}

\begin{warningbox}[Önemli yorum sınırı]
Repo tarihinde ALFA üzerinde mevcut Isolation Forest skoru ile CUSUM karar
katmanının birleştiği bir hücre, advisory Gate-C benzeri hedefi
\metric{0,625 recall / 7,91 \fasaat} ile geçmiştir. Bu sonuç:
\begin{itemize}
  \item yeni LightGBM modelinin başarısı değildir;
  \item çok küçük ve platforma özel ALFA havuzuna aittir;
  \item bağımsız final kör holdout ile doğrulanmamıştır;
  \item diğer platformlara ve genel ürün iddiasına taşınamamıştır.
\end{itemize}
Bu nedenle raporun doğru cümlesi ``hiçbir sayı hiçbir zaman geçmedi'' değil,
``geçerli event truth, gerçekçi alarm yükü, platform genellemesi ve bağımsız
final doğrulamayı birlikte sağlayan genel bir uçtan uca aday oluşmadı''dır.
\end{warningbox}

\tableofcontents
\clearpage

\pagenumbering{arabic}

\chapter{Yönetici Özeti}

\section{Ne yapıldı?}

Çalışmada beş ana veri ailesi kullanıldı:

\begin{enumerate}
  \item \textbf{ALFA}: sabit kanatlı İHA gerçek aktüatör ve motor arızaları,
  \item \textbf{UAV Attack}: PX4 multirotor GPS spoofing, jamming ve ping/DoS kayıtları,
  \item \textbf{UAV-SEAD}: daha geniş PX4 normal ve anomali havuzu,
  \item \textbf{RflyMAD}: gerçek motor ve sensör arızaları ile simülasyon adayları,
  \item \textbf{ADS-B/adsb.lol}: etiketsiz, gerçek ve büyük ölçekli insanlı hava aracı trafiği.
\end{enumerate}

Bu kaynaklar üzerinde:

\begin{itemize}
  \item parser, zaman hizalama, birim dönüşümü ve null/missingness yönetimi;
  \item fiziksel tutarlılık residual'ları ve nedensel zaman serisi özellikleri;
  \item gerçek event interval'ı ve sentetik enjeksiyon truth sözleşmeleri;
  \item Isolation Forest, LightGBM, Dense-AE, LSTM-AE, LSTM forecasting,
  USAD ve Chronos;
  \item robust median/MAD, rolling istatistik, EWMA, CUSUM, persistence,
  K-of-N ve episode birleştirme;
  \item platforma özgü eşik kalibrasyonu ve conformal calibration;
  \item satır, pencere, event, uçuş ve uçuş-saati seviyesinde değerlendirme
\end{itemize}

uygulandı.

\section{Neden tek bir başarı metriği yeterli olmadı?}

İlk deneylerde AUROC veya AUPRC yüksek olduğunda model başarılı görünebiliyordu.
Fakat operasyonel anomali tespitinde asıl soru, anomalilerin ortalamada daha
yüksek skorlanması değil; kabul edilebilir alarm bütçesinde kaç gerçek olayın
yakalandığıdır.

Örneğin ADS-B robust rule skorlayıcı, corrected truth v2 üzerinde
\metric{0,764883 AUROC} ve \metric{0,883313 AUPRC} üretmiştir. Aynı eşik doğal
trafikte yaklaşık \metric{4,808533 alarm olayı/saat} oluşturmuş ve uçuşların
yaklaşık \metric{\%89,24}'ünde alarm vermiştir. Dolayısıyla ranking metriği iyi,
ürün çalışma noktası ise kabul edilemezdir.

\section{Nihai deneysel örüntü}

Tekrarlanan sonuç şudur:

\begin{center}
\begin{tcolorbox}[
  width=0.92\textwidth,
  colback=LightBlue,
  colframe=Blue,
  boxrule=0.7pt,
  arc=2mm
]
\centering\sffamily
\textbf{Eşik gevşetildiğinde} recall yükseldi fakat yanlış alarm patladı.\\[0.35em]
\(\Downarrow\)\\[-0.2em]
\textbf{Eşik operasyonel bütçeye sıkılaştırıldığında} yanlış alarm azaldı fakat recall sıfıra yaklaştı.
\end{tcolorbox}
\end{center}

Model aileleri ve feature setleri değişse de bu değiş-tokuş genel olarak
kapanmadı. Sorun yalnızca tek bir kötü algoritma, yanlış hiperparametre veya
yetersiz eğitim süresi değildir.

\section{Yönetim için önerilen kapanış cümlesi}

\begin{infobox}[Önerilen resmî ifade]
Projede fiziksel kural tabanlı residual analizleri, robust istatistiksel
skorlar, EWMA/CUSUM, Isolation Forest, LightGBM, Dense ve LSTM Autoencoder,
LSTM forecasting, USAD, Chronos ve contextual-conformal modeller test
edilmiştir. Veri genişletme, platforma özgü kalibrasyon, session tabanlı
split, gerçek arıza interval'ı, sentetik severity enjeksiyonu, nedensel
özellik üretimi, missingness/staleness modelleme ve operasyonel alarm bütçesi
değerlendirmeleri uygulanmıştır. Bazı dar veri setlerinde veya sentetik
anomali türlerinde iyi ranking sonuçları alınmış olsa da bu sonuçlar gerçek
event truth, platform genellemesi, düşük yanlış alarm ve bağımsız final
doğrulamayı birlikte karşılamamıştır. Mevcut veri ve sensör kapsamı altında
genel amaçlı operasyonel anomali dedektörü hedefi teknik olarak achievable
değildir.
\end{infobox}

\chapter{Kapsam, Deney Sözleşmesi ve Metrikler}

\section{Repo kapsamı}

Rapor iki tarihsel alanı ve aktif alanı birlikte ele alır:

\begin{itemize}
  \item \code{archive/2026-07-10_legacy_non_adsb_ml}: ALFA, UAV Attack,
  UAV-SEAD ve RflyMAD üzerinde ML-0--ML-16 araştırma hattı;
  \item \code{archive/2026-07-10_rejected_adsb_attempts}: ilk iki,
  koordinesiz ve daha sonra reddedilmiş ADS-B yaklaşımı;
  \item aktif \code{adsb/}: 10 Temmuz resetinden sonra temiz namespace ile
  kurulan fiziksel residual, kural, CUSUM, S2 ve contextual model hattı.
\end{itemize}

Arşiv tarihçedir; aktif modele arşiv paketlerinden kod import edilmemiştir.
Aktif ADS-B parser'ı veri okuma altyapısıdır; tek başına kabul edilmiş bir
anomali modeli değildir.

\section{Veri katmanları}

\begin{longtable}{L{0.16\textwidth}L{0.30\textwidth}L{0.45\textwidth}}
\toprule
\textbf{Katman} & \textbf{İşlev} & \textbf{Temel kural}\\
\midrule
\endfirsthead
\toprule
\textbf{Katman} & \textbf{İşlev} & \textbf{Temel kural}\\
\midrule
\endhead
Bronze & Ham kaynağın saklanması & Kaynak dosyası değiştirilmez; provenance ve indirme bilgisi korunur.\\
Silver & Kaynak-özel ayrıştırma & Birimler, topic'ler, zaman ekseni, veri erişilebilirliği ve kaynak semantiği düzenlenir.\\
Gold/Feature & Model girdileri & Fiziksel residual, rolling, CUSUM, spectral, missingness ve platforma özgü özellikler üretilir.\\
Split & Train/calibration/test/holdout & Uçuş veya session izolasyonu uygulanır; normal-only novelty training ile supervised split birbirinden ayrılır.\\
Artifact & Model, scaler ve eşik & Eğitim verisine ait scaler, CUSUM baseline, model ağırlığı, karar politikası ve checksum birlikte tutulur.\\
\bottomrule
\end{longtable}

Rapor tarihi civarındaki yerel envanter yaklaşık \(41{,}6\ \text{GB}\) veri
kapsamına ulaşmıştır; Bronze, Silver ve sentetik ADS-B çıktıları bunun büyük
bölümünü oluşturmuştur. Bu hacim, bağımsız arıza örneği sayısının da büyük
olduğu anlamına gelmez.

\section{Değerlendirme birimleri}

\begin{longtable}{L{0.17\textwidth}L{0.34\textwidth}L{0.40\textwidth}}
\toprule
\textbf{Birim} & \textbf{Ne ölçer?} & \textbf{Yanlış kullanım riski}\\
\midrule
\endfirsthead
\toprule
\textbf{Birim} & \textbf{Ne ölçer?} & \textbf{Yanlış kullanım riski}\\
\midrule
\endhead
Satır & Tek telemetri satırının anomalilik skoru & Aynı uçuş içindeki binlerce bağımlı satır bağımsız örnek sanılabilir.\\
Pencere & Kısa zaman dizisinin skoru & Tüm saldırı uçuşunu pozitif etiketlemek kısa imzayı seyreltebilir.\\
Event onset & Anomali başladıktan sonra yeni alarm oluşması & Event başlamadan önce açık alarm ``tespit'' sayılırsa recall şişer.\\
Event overlap & Alarm ile anomali aralığının herhangi bir kesişimi & Onset duyarlılığı ve gerçek yeni tespit gücüyle aynı değildir.\\
Uçuş & Uçuşta en az bir alarm veya uçuş agregat skoru & Uzun uçuşların alarm alma olasılığı doğal olarak yükselir.\\
Uçuş-saati & Zamanla normalize edilmiş alarm yükü & Episode merge alarm saturasyonunu gizleyebilir; row coverage ile birlikte bakılmalıdır.\\
\bottomrule
\end{longtable}

\section{Başarı kapıları}

\begin{longtable}{L{0.18\textwidth}L{0.72\textwidth}}
\toprule
\textbf{Kapı} & \textbf{Tanım}\\
\midrule
\endfirsthead
\toprule
\textbf{Kapı} & \textbf{Tanım}\\
\midrule
\endhead
Gate A & Determinizm, veri sızıntısı olmaması, split ve holdout izolasyonu, checksum ve güvenlik doğrulamaları.\\
Gate B & Yeni adayın kayıtlı mevcut baseline'ı önceden belirlenmiş marj ve seed kararlılığıyla geçmesi. Bazı fazlarda \(+0{,}05\) recall ve en az 3/5 seed şartı kullanıldı.\\
Gate C -- critical & En az \(0{,}30\) event recall ve en fazla \(2\) \fasaat.\\
Gate C -- advisory & En az \(0{,}50\) event recall ve en fazla \(12\) \fasaat.\\
ADS-B Pareto & 100 skorlanabilir uçuş-saatinde \(0{,}1,\ 0{,}5,\ 1,\ 2,\ 5\) episode bütçeleri.\\
\bottomrule
\end{longtable}

\begin{warningbox}[Kapılar arasında doğrudan sayı karşılaştırması]
Eski Gate-C birimi \(\FA/\text{saat}\), aktif ADS-B Pareto birimi ise
\(\text{episode}/100\text{ saat}\)tir. Bu iki bütçe aynı değildir. ADS-B
bütçesi çok daha sıkıdır ve ürün sözleşmesinin değiştiğini gösterir. Bu
nedenle eski bir recall/FA hücresi ile ADS-B contextual sonucu doğrudan
``hangisi daha iyi?'' diye kıyaslanmamalıdır.
\end{warningbox}

\chapter{Veri Setleri, Temizlik ve Feature Engineering}

\section{ALFA}

\subsection{Kaynak ve etiket yapısı}

ALFA, sabit kanatlı bir İHA'da gerçek motor ve kumanda yüzeyi arızaları
içeren etiketli uçuş veri setidir. Resmî işlenmiş corpus 47 uçuştu. Ham ROS
bag kaynaklarından 7 ek uçuş kurtarılarak legacy feature havuzu 54 uçuşa
çıkarıldı.

Sınıflar normal, engine, aileron, rudder, elevator, birleşik
aileron--rudder ve unknown kayıtlarını içerir. Bazı sınıflarda yalnızca
birkaç bağımsız uçuş vardır. Repo içindeki farklı artifact revizyonlarında
normal veya rudder sayısı bir uçuş farklı raporlanabilmektedir; bu nedenle
tek tek artifact manifestleri karıştırılmadan kullanılmalıdır.

\begin{warningbox}[İstatistiksel örneklem]
ALFA'nın milyonlarca telemetri satırı üretmesi, rudder veya elevator sınıfı
için büyük bir örneklem oluşturmaz. Bağımsız birim uçuş/session olduğundan
bazı arıza sınıflarındaki recall değerleri bir veya birkaç örneğe dayanır.
\end{warningbox}

\subsection{Preprocessing ve temizlik}

\begin{itemize}
  \item Global position tablosu zaman omurgası olarak kullanıldı.
  \item Topic CSV'leri en yakın zaman damgasıyla \code{merge_asof}
  kullanılarak, aktif parser'da yaklaşık \(0{,}5\) saniye toleransla birleştirildi.
  \item Roll, pitch, yaw ve airspeed için measured--commanded farkları üretildi.
  \item \code{nav_info} üzerinden altitude error, airspeed error,
  cross-track error ve waypoint distance alındı.
  \item \code{mavctrl} kaynaklarından \(x/y/z\) path deviation eklendi.
  \item \code{vfr_hud} üzerinden airspeed, ground speed, throttle ve climb eklendi.
  \item Ölçülen ve istenen \(x/y/z\) hızları ile vektörel hız büyüklükleri türetildi.
  \item Mutlak enlem, boylam ve mutlak zaman doğrudan ML girdisi yapılmadı.
  \item Arıza tipi klasör adından, aktif arıza başlangıcı mümkün olduğunda
  \code{failure_status} üzerinden çıkarıldı.
  \item Çoğul/tekil etiket farkları normalize edildi.
\end{itemize}

Eski referans parser daha agresif bir politika kullanıyordu: tümü-null
kolonları atıyor, yaklaşık \(250\) ms toleransla birleştiriyor ve sayısal
kolonlara lineer interpolasyon uyguluyordu. Bu davranış aktif parser ile
aynı kabul edilmemelidir.

\subsection{Bulunan veri hatası}

\code{velocity_mps} ilk üretimlerde \(\%100\) null çıkıyordu. Parser,
\code{measured} ve \code{desired} benzeri adlar beklerken gerçek alanlar
\code{field.meas_x} ve \code{field.des_x} biçimindeydi. Kolon eşleme
düzeltildikten sonra null oranı yaklaşık \(\%9{,}3\)'e indi.

\subsection{Üretilen özellikler}

\begin{itemize}
  \item roll, pitch, yaw, airspeed ve velocity residual'ları;
  \item residual mutlak değeri ve zamansal türevi;
  \item GPS haversine hızı, ivme ve iki nokta arası mesafe;
  \item altitude, airspeed, cross-track ve path-deviation residual'ları;
  \item 2 ve 5 saniyelik rolling ortalama, standart sapma ve enerji;
  \item EWMA, nedensel CUSUM ve spektral enerji;
  \item freeze sayaçları, availability maskeleri ve \code{in_air};
  \item koordine dönüş residual'ı:
  \[
    r_{\text{turn}} =
    \dot{\psi} - \frac{g\tan(\phi)}{V},
    \qquad V>5\ \text{m/s};
  \]
  \item özgül enerji ve tırmanma tutarlılığı.
\end{itemize}

\subsection{Sentetik enjeksiyonlar}

ALFA Silver verisine yalnız kontrollü değerlendirme amacıyla freeze, bias,
drift, noise, GPS ramp ve dropout enjeksiyonları uygulandı. Enjeksiyonlar
gerçek arıza ground-truth'unun yerine geçirilmedi.

\subsection{ALFA deney sonuçları}

\begin{longtable}{L{0.34\textwidth}L{0.22\textwidth}L{0.35\textwidth}}
\toprule
\textbf{Deney} & \textbf{Sonuç} & \textbf{Yorum}\\
\midrule
\endfirsthead
\toprule
\textbf{Deney} & \textbf{Sonuç} & \textbf{Yorum}\\
\midrule
\endhead
Monolitik satır Isolation Forest & ROC-AUC \(0{,}497\pm0{,}026\) & Rastgele seviyesinde; satır etiketi ve feature sulanması sorunu.\\
Modüler Isolation Forest & Uçuş ROC-AUC \(0{,}833\pm0{,}172\) & Modüler fizik grupları monolitik modele göre daha iyi.\\
Guidance-only IF & ROC-AUC yaklaşık \(0{,}864\) & Kontrol-tepki modülü bazı uçuşlarda gürültü ekledi.\\
İlk LSTM-AE & ROC-AUC yaklaşık \(0{,}731\) & Normal eğitim verisi yetersizdi.\\
Dense-AE & ROC-AUC yaklaşık \(0{,}622\) & LSTM ve IF altında kaldı.\\
Normal veri genişletilmiş LSTM-AE & \(0{,}918\pm0{,}104\) & Projedeki gerçek DL kazanımlarından biri; ancak küçük \(n\) ve geniş belirsizlik.\\
Causal düzeltme öncesi \code{alt_error_cusum} & \(0{,}878\) & Gelecek bilgisi/seri-geneli baseline şüphesi.\\
Causal düzeltme sonrası \code{alt_error_cusum} & \(0{,}611\) & Eski görünen başarı büyük ölçüde düştü.\\
Causal en güçlü tek sinyal & \code{xtrack_error}: \(0{,}751\) & Daha savunulabilir residual.\\
Sabit reçete LightGBM / IF / LSTM-AE & \(0{,}843/0{,}858/0{,}872\) AUPRC & Supervised LightGBM mevcut yöntemleri geçmedi.\\
Mevcut IF + CUSUM advisory & Recall \(0{,}625\), FA \(7{,}91/\text{saat}\) & Dar ALFA hücresi geçti; genel ürün ve yeni model başarısı değil.\\
\bottomrule
\end{longtable}

\section{UAV Attack}

\subsection{Kaynak}

PX4 multirotor kayıtlarından oluşan 19 log işlendi:

\begin{itemize}
  \item 6 benign,
  \item 6 GPS spoofing,
  \item 6 ping/DoS,
  \item 1 jamming.
\end{itemize}

Kaynak yaklaşık \(683{,}9\) MB ve 767 CSV dosyası içeriyordu.

\subsection{Preprocessing}

Aktif legacy parser dört ana topic kullandı:

\begin{itemize}
  \item \code{vehicle_global_position},
  \item \code{vehicle_attitude},
  \item \code{battery_status},
  \item \code{vehicle_gps_position}.
\end{itemize}

Dosya adı regex'i düzeltildi. Etiket çıkarımında en yakın klasöre öncelik
verildi; böylece üst klasöründe ``GPS Spoofing'' yazan benign logların
yanlış etiketlenmesi önlendi. Ping/DoS için ayrıca kural eklendi.

Global position zaman omurgası olarak kullanıldı; quaternion Euler
açılarına çevrildi. Attitude, battery ve GPS alanları yaklaşık \(200\) ms
toleransla birleştirildi. NED sistemindeki aşağı-pozitif hız, yukarı-pozitif
vertical-rate biçimine çevrildi. Varsa UTC mikro-saniye alanı, yoksa göreli
zaman kullanıldı.

Batarya akımındaki \(-1\) değerinin fiziksel akım değil ``bilinmiyor''
sentineli olduğu görülerek maskelendi.

\subsection{Özellikler}

\begin{itemize}
  \item Haversine ve Öklid tabanlı kinematik;
  \item bildirilen GPS hızı ile konum türevi arasındaki
  \code{gps_speed_residual};
  \item dikey hız residual'ı;
  \item course change;
  \item GPS health ve uydu kalitesi değişimleri;
  \item batarya özellikleri;
  \item rolling, CUSUM, freeze ve staleness;
  \item missingness ve \code{is_moving}.
\end{itemize}

\subsection{Sonuç ve gözlenebilirlik sınırı}

\begin{longtable}{L{0.36\textwidth}L{0.23\textwidth}L{0.32\textwidth}}
\toprule
\textbf{Deney} & \textbf{Sonuç} & \textbf{Yorum}\\
\midrule
\endfirsthead
\toprule
\textbf{Deney} & \textbf{Sonuç} & \textbf{Yorum}\\
\midrule
\endhead
Monolitik satır IF & \(0{,}209\pm0{,}138\) ROC-AUC & Skor ters yönde; NaN-medyan imputasyonu saldırı satırını normale çekti.\\
Modüler IF & \(0{,}600\pm0{,}212\) uçuş ROC-AUC & Monolitik yaklaşımdan iyi, belirsizlik yüksek.\\
LSTM/Dense AE & Yaklaşık \(0{,}677\) & Spoofing/jamming güçlü; ping/DoS zayıf.\\
LSTM-AE spoofing/jamming & Yaklaşık \(1{,}0/1{,}0\) & Mevcut sensörlerde güçlü imza olan sınıflar.\\
LSTM-AE ping/DoS & Yaklaşık \(0{,}53\) & Rastgele seviyesine yakın.\\
SEAD eşiği transferi ve yeniden kalibrasyon & \(0{,}375\to0{,}783\); normal FA \(1{,}00\to0{,}33\) & Eşiğin platforma özel olması gerektiğini gösterdi.\\
\bottomrule
\end{longtable}

Ping/DoS loglarının 4/6'sında kullanılan topic'lerde gözlenebilir saldırı
etkisi bulunmadı. Paket geliş aralığı veya ağ katmanı metadatası mevcut
değildi. Modelin görmediği bir ağ etkisini uçuş telemetrisinden tespit etmesi
beklenemez; bu algoritma değil sensör/gözlenebilirlik sınırıdır.

2 m/s düşük şiddetli GPS ramp enjeksiyonlarında uçuş alarmı oluşmadı. Daha
geniş 2--20 m/s severity sweep'i tamamlanmadığı için detectability eşiği
nicel olarak belirlenemedi.

\section{UAV-SEAD}

\subsection{Havuz ve split}

Güncel feature manifesti toplam 1.244 uçuş içerir:

\begin{center}
\begin{tabular}{lr}
\toprule
\textbf{Sınıf} & \textbf{Uçuş}\\
\midrule
Normal & 898\\
Altitude anomaly & 72\\
External position anomaly & 193\\
Global position anomaly & 40\\
Mechanical anomaly & 41\\
\midrule
Toplam & 1.244\\
\bottomrule
\end{tabular}
\end{center}

200 uçuş final holdout olarak ayrılmış, geliştirme havuzu 1.044 uçuşta
kalmıştır. Eski belgelerdeki ``1.044 toplam'' ifadesi tüm yerel havuz değil,
holdout dışı geliştirme kapsamıdır. Veri 60, 179, 349, 611 ve 1.244 uçuş
aşamalarıyla genişletildi.

Mapping içindeki sekiz multi-class kategori ve 141 ``Uncategorized'' kayıt,
tekil ve doğrulanabilir sınıf sözleşmesi sağlamadığı için resmi havuza
alınmadı.

\subsection{Preprocessing}

\code{pyulog} ile aşağıdaki PX4 topic aileleri işlendi:

\begin{itemize}
  \item global/local position ve attitude;
  \item battery ve GPS;
  \item estimator status ve EKF innovations;
  \item barometer;
  \item setpoint;
  \item actuator controls ve actuator outputs.
\end{itemize}

Indoor uçuşlarda global position bulunmadığında local position fallback
olarak kullanıldı. Topic'ler yaklaşık \(200\) ms toleransla birleştirildi.
Mapping aralıkları event truth ve kategori üretimi için korundu.

Split session tarihleri üzerinden yapıldı. Aynı session'ın train ve test
tarafına düşmesi engellendi; normal ve anomali karışık session'lar
karantinaya alındı.

\subsection{Eklenen feature aileleri}

\begin{itemize}
  \item altitude--barometer--local position residual'ları;
  \item EKF innovation büyüklükleri, oranları ve status flag'leri;
  \item setpoint response residual'ları;
  \item actuator effort ve control strain;
  \item nedensel/expanding actuator output symmetry;
  \item vibration, reset counter ve valid maskeleri;
  \item dikey innovation bileşenleri.
\end{itemize}

\subsection{Yapısal veri problemleri}

\begin{itemize}
  \item \code{alt_local_residual}, altitude anomaly grubunda \(\%0\)
  erişilebilirlik gösterdi.
  \item Barometer kapsaması yaklaşık \(\%7\) düzeyindeydi.
  \item EKF oranları bazı testlerde ters ayrım üretti; estimator ölçümü
  reddettiğinde innovation bastırılabildiği için anomali daha normal
  görünebiliyordu.
  \item Yüzlerce normal uçuş yalnızca onlarca bağımsız session'a aitti.
  Örneğin 398 normal örnek 64 session'dan, geliştirme tarafı yaklaşık 49
  session'dan geliyordu.
\end{itemize}

\subsection{UAV-SEAD sonuçları}

\begin{longtable}{L{0.37\textwidth}L{0.23\textwidth}L{0.31\textwidth}}
\toprule
\textbf{Faz/aday} & \textbf{Sonuç} & \textbf{Karar}\\
\midrule
\endfirsthead
\toprule
\textbf{Faz/aday} & \textbf{Sonuç} & \textbf{Karar}\\
\midrule
\endhead
Naif range satır ROC & \(0{,}474\to0{,}799\) & Artışın önemli kısmı veri/split metodolojisinden geldi.\\
LightGBM descriptor & AUPRC \(0{,}349\) & IF \(0{,}385\)'in altında; Gate B kaldı.\\
ML9 fusion & Recall \(0{,}222\), FA \(25{,}83/\text{saat}\) & Operasyonel Gate C kaldı.\\
Chronos motor & \(0{,}205\to0{,}390\) recall & Mechanical Gate B geçti.\\
Chronos vertical & \(0{,}096\to0{,}023\) recall & Kötüleşti.\\
Chronos fusion & \(0{,}213/23{,}92\) & Kategori kazancı genel alarma dönüşmedi.\\
Tek-feature \code{actuator_thrust_cmd} & \(0{,}205\to0{,}459\) & Gate B geçti; normal yük \(38{,}1\) FA/saat.\\
İki-kanal dengeli birleşim & Recall \(0{,}217\to0{,}291\), FA \(23{,}7\to44{,}7\) & Recall, alarm yüküyle satın alındı; reddedildi.\\
ML14 normal genişletme & Existing fusion FA \(23{,}63\to9{,}86\) & FA düştü; recall \(0{,}1186\)'ya çöktü.\\
ML14 critical & FA \(12{,}87\to1{,}60\), recall \(0{,}0417\) & Bütçe içinde fakat faydasız recall.\\
ML15 drift/jackknife en iyi & \(0{,}1145/8{,}4146\) & Gate C/D3 kaldı.\\
ML15 critical & \(0{,}0433/1{,}596\) & Alarm bütçesi sağlanırken gerçek tespit kayboldu.\\
\bottomrule
\end{longtable}

\subsection{Derin öğrenmede magnitude artefaktı}

LSTM-AE, Dense-AE ve USAD skorları ilk bakışta yaklaşık \(0{,}21\) recall
gösterebiliyordu. Fakat:

\begin{itemize}
  \item eğitimli ve eğitimsiz model skor korelasyonu yaklaşık \(0{,}964\),
  \item eğitimli skor ile ham residual magnitude korelasyonu yaklaşık \(0{,}965\)
\end{itemize}

ölçüldü. Bu, skorun büyük ölçüde öğrenilmiş davranıştan değil genlikten
geldiğini gösterdi. Relative error ile magnitude korelasyonu
\(0{,}15\)--\(0{,}55\) seviyesine düşürüldüğünde recall çoğu deneyde
\(\%6\)'nın altında kaldı.

\section{RflyMAD}

\subsection{Gerçek uçuş havuzu}

\begin{center}
\begin{tabular}{lr}
\toprule
\textbf{Kategori} & \textbf{Uçuş}\\
\midrule
Real-Motor & 242\\
Real-Sensors & 197\\
Real-No\_Fault & 51\\
\midrule
Toplam & 490\\
\bottomrule
\end{tabular}
\end{center}

SEAD PX4 parser altyapısı yeniden kullanıldı. Klasör taksonomisinden sınıf,
case kimliğinden session üretildi.

\subsection{Event truth düzeltmesi}

Başlangıçta klasörü arızalı olan uçuşun tamamı anomali kabul edilmişti.
Gerçek arıza aralığı daha sonra:

\begin{itemize}
  \item öncelikle \code{rfly_ctrl_lxl} içinde \code{id/mode != 1500},
  \item mümkün değilse \code{TestInfo}
\end{itemize}

üzerinden çıkarıldı. Klasörü motor fault olduğu halde aktif fault aralığı
bulunmayan beş uçuş corrected scoring'den dışlandı.

Pooled kaynaklarda bir platformda hiç bulunmayan kolonun başka platform
medyanıyla doldurulması da engellendi. Source-schema presence maskesi
eklenerek ``hayalet sensör'' imputasyonu düzeltildi.

\subsection{İlk görünen başarı ve geçersizleşmesi}

\begin{longtable}{L{0.40\textwidth}L{0.23\textwidth}L{0.28\textwidth}}
\toprule
\textbf{Değerlendirme} & \textbf{Sonuç} & \textbf{Yorum}\\
\midrule
\endfirsthead
\toprule
\textbf{Değerlendirme} & \textbf{Sonuç} & \textbf{Yorum}\\
\midrule
\endhead
Eski full-flight proxy advisory & \(0{,}749/10{,}18\) & İlk bakışta geçiyor; gerçek event interval kullanılmadı.\\
Eski full-flight proxy critical & \(0{,}573/1{,}23\) & Geçersiz proxy başarı.\\
Corrected RFLY-only \code{itki_komutu} CUSUM advisory & \(0{,}526/22{,}275\) & Recall kısmen korundu, alarm yükü başarısız.\\
Corrected critical & \(0{,}442/9{,}228\) & Critical bütçeyi aştı.\\
Pooled corrected en iyi & Yaklaşık \(0{,}149/30\) & Platform havuzlama kötüleşti.\\
Bütçe içine zorlanan pooled & Yaklaşık \(0{,}066/9{,}39\) & Recall operasyonel anlamını kaybetti.\\
\bottomrule
\end{longtable}

Simülasyon ve severity çalışması planlandı; resmi deney yürütüldüğü sırada
gerekli simülasyon/\code{TestInfo} kaynakları yerelde olmadığı için preflight
aşamasında durdu. Simülasyon subsetlerinin daha sonra indirilmiş olması,
tamamlanmış ve geçerli bir deney sonucu anlamına gelmez.

\section{Fiziksel tutarsızlık nedeniyle eğitim dışı bırakılan ek CSV'ler}

SurveilDrone-Net23 benzeri bir drone kaynağı ile afet operasyonları temalı
başka bir CSV adayı incelendi. GPS teleportation, yalnız tamsayı irtifa,
konum--hız tutarsızlığı ve fiziksel olarak mümkün görünmeyen hareketler
bulundu. Bu kaynaklar gerçek anomaly ground-truth'u olarak kullanılmadı;
yalnız pipeline veya dashboard demosu için uygun görüldü.

\chapter{ADS-B Veri Hattı ve Fiziksel Model}

\section{Reset sonrası problem çerçevesi}

10 Temmuz 2026 tarihinde eski non-ADS-B hat ile ilk iki ADS-B yaklaşımı
arşivlendi. Aktif çalışma temiz \code{adsb/} namespace'iyle başladı.

ADS-B kaynağı gerçek, büyük ölçekli fakat etiketsiz insanlı hava aracı
trafiğidir. Bu nedenle aktif problem supervised saldırı sınıflandırması
değil, fiziksel/bağlamsal beklenmediklik ve veri kalitesi izleme problemidir.

\section{Veri envanteri}

Üç tarihsel gün kullanıldı:

\begin{itemize}
  \item 28 Şubat 2026,
  \item 1 Mart 2026,
  \item 16 Mart 2026.
\end{itemize}

Authoritative streaming envanteri:

\begin{itemize}
  \item 638 Silver part,
  \item 256.155.009 satır,
  \item yaklaşık 542.461 uçuş segmenti,
  \item yaklaşık 497.570,9 skorlanabilir uçuş-saat.
\end{itemize}

Eski dokümanlarda 256.150.550 satır sayısı da bulunur. Manifestler arasında
4.459 satırlık sayım farkı kayıtlıdır; raporda streaming authoritative sayı
esas alınmıştır.

\subsection{Alan kapsaması}

\begin{center}
\begin{tabular}{lr@{\qquad}lr}
\toprule
\textbf{Alan} & \textbf{Kapsama} & \textbf{Alan} & \textbf{Kapsama}\\
\midrule
Latitude/longitude & \(\%100\) & Speed & \(\%98{,}1\)\\
Track & \(\%95{,}2\) & Barometric altitude & \(\%89{,}4\)\\
Geometric altitude & \(\%89{,}2\) & Vertical rate & \(\%89{,}4\)\\
IAS & \(\%29{,}7\) & Roll & \(\%28{,}4\)\\
\bottomrule
\end{tabular}
\end{center}

Doğrulanmış UAV/drone etiketi bulunmamaktadır. B6 benzeri 25 örnek
incelenmiş fakat UAV kimliği doğrulanmamıştır. Dolayısıyla aktif veri,
doğrudan ``drone saldırı veri seti'' değildir.

\section{Parser ve veri semantiği}

\begin{itemize}
  \item Gzip veya plain trace JSON okundu.
  \item 14 elemanlı trace satırları ayrıştırıldı.
  \item Feet metreye \(0{,}3048\), knot m/s'ye \(0{,}5144\), feet/min m/s'ye
  \(0{,}00508\) katsayılarıyla çevrildi.
  \item Doğrudan trace alanlarına forward-fill uygulanmadı.
  \item Provider'ın sparse aircraft-dictionary alanları kendi güncelleme
  semantiğiyle taşındı; v2'de her alan için update flag, update time ve age
  kaydedildi.
  \item \code{on_ground} olduğunda altitude null yapıldı.
  \item \code{dbFlags} military biti ayrıştırıldı; alan eksikse ``military
  değil'' değil ``doğrulanmadı'' kabul edildi.
  \item Part üretimi streaming/chunked yapıldı.
  \item Checkpoint ve resume atomik hale getirildi.
  \item Yarım kalan tar arşivinde yalnız ilgili kısmi part'lar kaldırılıp o
  tar yeniden işlendi.
\end{itemize}

Uçuş segmentasyonu 1.800 saniyelik boşluk kuralıyla yapıldı. 1.500 aircraft
üzerindeki kontrol örneği 4.230 uçuş üretmiş, provider
\code{flags_new_leg} ile uyum yaklaşık \(\%60{,}4\) kalmıştır. Uçuş sınırı
kendi başına kesin ground truth değildir.

\section{Fiziksel residual'lar}

\subsection{Vertical-rate residual}

\[
r_{vr,t}
= vr_{\text{reported},t}
- \frac{alt_t-alt_{t-1}}{\Delta t}.
\]

\subsection{Ground-speed residual}

\[
r_{v,t}
= v_{\text{reported},t}
- \frac{d_{\text{haversine}}(p_t,p_{t-1})}{\Delta t}.
\]

\subsection{Heading residual}

\[
r_{\psi,t}
= \operatorname{wrap}_{[-180,180)}
\left(
\psi_{\text{reported},t}
- \operatorname{bearing}(p_{t-1},p_t)
\right).
\]

\subsection{East/north hız residual'ları}

Bildirilen ground speed ve track üzerinden doğu/kuzey hız bileşenleri
çıkarıldı; bunlar konum türevinden elde edilen bileşenlerle karşılaştırıldı:

\[
r_{E,t}=v^{\text{reported}}_{E,t}-v^{\text{position}}_{E,t},
\qquad
r_{N,t}=v^{\text{reported}}_{N,t}-v^{\text{position}}_{N,t}.
\]

\subsection{Diğer fiziksel kontroller}

\begin{itemize}
  \item Barometric--geometric altitude farkının türevi,
  \item gözlenen track-rate ile \(g\tan(\text{roll})/V\) arasındaki
  turn--bank residual'ı,
  \item cadence, message gap ve availability maskeleri.
\end{itemize}

\section{S2 veri kalitesi katmanı}

Hareket residual'larından ayrı ikinci katmanda squawk, emergency, NIC,
NACp, SIL, message gap, altitude availability ve veri kalitesi izlendi.

\begin{center}
\begin{tabular}{lr}
\toprule
\textbf{Doğal kalite nedeni} & \textbf{Yaklaşık olay/saat}\\
\midrule
Message gap & 2,8905\\
NIC unknown & 0,5281\\
NACp missing & 0,1952\\
SIL missing & 0,1952\\
Altitude unavailable & 0,00508\\
Declared emergency/status & 0,00033\\
\bottomrule
\end{tabular}
\end{center}

Bu değerler attack ground-truth'u değildir; doğal veri kalitesi ve beyan
durumu yükünü gösterir.

\section{Sentetik enjeksiyon truth v2}

8.910 uçuş için clean ve beş kontrollü enjeksiyon tarifi üretildi:

\begin{enumerate}
  \item vertical-rate freeze,
  \item ground-speed bias,
  \item track freeze,
  \item 2 m/s position ramp,
  \item altitude dropout.
\end{enumerate}

Toplam yaklaşık 26,8 milyon satır ve 646 MB çıktı oluştu.

Truth v2 üç farklı kavramı ayırdı:

\begin{itemize}
  \item \code{injection_active}: kod enjeksiyonu o anda aktif mi?
  \item \code{observable_changed}: modelin kullandığı gözlem gerçekten değişti mi?
  \item \code{evaluable_truth}: satır değerlendirme için fiziksel ve
  sözleşmesel olarak uygun mu?
\end{itemize}

Eski proxy etiket, enjeksiyonun aktif olduğu her satırı pozitif sayarak
yaklaşık \(0{,}75\) civarında yapay AUC tavanı üretiyordu. Truth v2 bu
yanlılığı azalttı. Sentetik veri yalnız kontrollü test içindir; gerçek saldırı
ground-truth'u olarak kabul edilmemiştir.

\chapter{Denenen Yöntemler ve Skorlar}

\section{Robust fiziksel kural skoru}

Her residual kanalı için eğitim-normal dağılımından median ve MAD hesaplandı:

\[
z_{c,t}
= \frac{|x_{c,t}-\operatorname{median}_c|}
{\operatorname{MAD}_c}.
\]

Kanal cezası:

\[
p_{c,t}
= \operatorname{clip}(z_{c,t}-3,0,10)
\]

olarak tanımlandı ve kullanılabilir kanallar eşit ağırlıkla birleştirildi.

İlk sürümde MAD değeri sıfır olan kanallara küçük epsilon verilmesi, normal
pencerelerin \(\%93{,}8\)'inin pozitif skor almasına neden oldu. Sıfır
MAD'li kanallar dışlanınca skor daha anlamlı hale geldi.

\begin{longtable}{L{0.36\textwidth}L{0.20\textwidth}L{0.35\textwidth}}
\toprule
\textbf{Truth-v2 sonucu} & \textbf{AUC/ölçüm} & \textbf{Yorum}\\
\midrule
\endfirsthead
\toprule
\textbf{Truth-v2 sonucu} & \textbf{AUC/ölçüm} & \textbf{Yorum}\\
\midrule
\endhead
Genel AUROC & 0,764883 & Ranking anlamlı.\\
Genel AUPRC & 0,883313 & Pozitif oranı ve reçete karışımı dikkate alınmalı.\\
Altitude dropout & 0,566965 & Fiziksel residual yerine S2 kapsamına daha uygun.\\
Speed bias & 0,974023 & Güçlü ve doğrudan gözlenebilir enjeksiyon.\\
Position ramp & 0,558927 & Zayıf/örtük gözlenebilirlik.\\
Track freeze & 0,889552 & İyi ayrım.\\
Vertical-rate freeze & 0,696337 & Orta düzey ayrım.\\
Doğal alarm yükü & 4,808533 episode/saat & Operasyonel olarak kabul edilemez.\\
Alarm alan uçuş oranı & \(\%89{,}2356\) & AUC başarısının ürün başarısı olmadığını gösterir.\\
\bottomrule
\end{longtable}

\section{CUSUM}

Kalıcı küçük sapmaları biriktirmek için tek yönlü CUSUM kullanıldı:

\[
S_t=\max\left(0,S_{t-1}+z_t-k\right).
\]

CUSUM uçuş sınırı, büyük zaman boşluğu ve zemin durumu gibi koşullarda reset
ediliyordu; fakat bir uçuş içinde uzun süre birikiyordu.

Full streaming değerlendirmede:

\begin{itemize}
  \item rule yükü yaklaşık \(1{,}061\)--\(1{,}171\) episode/saat,
  \item CUSUM \(h=1\) yükü yaklaşık \(5{,}33\)--\(6{,}07\) episode/saat,
  \item uçuşların yaklaşık \(\%99{,}1\)'i alarm,
  \item satırların yaklaşık \(\%78\)--\(\%80\)'i alarm
\end{itemize}

oldu. Yakın alarmları tek episode olarak birleştirmek gerçek saturasyonu
gizleyebildiği için row coverage ve alarm alan uçuş oranı da raporlandı.

\section{Isolation Forest}

Isolation Forest, etiketsiz normal davranıştan dışarı çıkan örnekleri bulmak
için kullanıldı. Monolitik, modüler, tek-feature ve flat contextual
varyantları denendi.

ADS-B'de aynı beş residual üzerinde:

\begin{itemize}
  \item eğitimli skor ile magnitude korelasyonu \(0{,}995958\),
  \item eğitimli skor ile shuffled-model skor korelasyonu \(0{,}995841\)
\end{itemize}

çıktı. Model pratikte residual genliğini yeniden sıralıyordu; bağlam
öğrendiği gösterilemedi.

\section{LightGBM}

10 saniyelik causal pencerelerden kanal başına yaklaşık 20 descriptor
üretilerek class-balanced LightGBM eğitildi. Normal-only novelty split'ten
ayrı supervised split ve session izolasyonu kullanıldı.

\begin{itemize}
  \item SEAD window AUPRC: LightGBM \(0{,}349\), IF \(0{,}385\);
  \item ALFA sabit reçete: LightGBM \(0{,}843\), IF \(0{,}858\),
  LSTM-AE \(0{,}872\).
\end{itemize}

Supervised model etiketi kullanmasına rağmen baseline'ı geçmedi. Optuna veya
sonuç sonrası hiperparametre kurtarma yapılmadı; ön-kayıtlı Gate B kararı
korundu.

\section{Dense-AE, LSTM-AE ve LSTM forecasting}

Autoencoder yaklaşımı normal pencereleri yeniden üretmeyi, reconstruction
error'ı anomaly skoru olarak kullanmayı amaçladı. LSTM forecasting geçmiş
pencereden bir sonraki davranışı tahmin etti.

ALFA'da normal veri artırıldığında LSTM-AE'nin \(0{,}731\to0{,}918\)
iyileşmesi, yeterli ve daha temsil edici normal verinin bazı dar domainlerde
faydalı olduğunu gösterdi.

ADS-B ilk DL turunda:

\begin{itemize}
  \item 60/638 part,
  \item window length 12,
  \item stride 6,
  \item history 8,
  \item maksimum gap 60 saniye,
  \item 15 epoch,
  \item batch size 512,
  \item sekiz ana residual kanalı,
  \item train-only robust scaling ve clip \(=5\)
\end{itemize}

kullanıldı.

Eski, truth hatası içeren pooled AUC'ler:

\begin{itemize}
  \item Dense-AE \(0{,}572\),
  \item LSTM-AE \(0{,}568\),
  \item LSTM forecast \(0{,}552\).
\end{itemize}

Skorların magnitude ile korelasyonu \(0{,}86\)--\(0{,}94\) aralığındaydı.

\section{USAD}

USAD iki decoder ve adversarial reconstruction ile daha keskin anomaly
ayrımı üretmek amacıyla denendi. Legacy küçük veri deneyinde ALFA \(0{,}450\),
UAV Attack \(0{,}531\) ile LSTM-AE'nin altında kaldı. ADS-B uygulamasında loss
yaklaşık 23 milyar seviyesine çıkarak sayısal olarak patladı.

Bounded decoder veya farklı optimizer çalışması araştırıldı fakat ana
fizik/truth hattı başarısızken USAD debug'ı ürün sonucunu değiştirecek öncelik
olarak görülmedi.

\section{Chronos ve foundation model denemeleri}

Chronos Bolt Tiny, SEAD'de \code{alt} ve actuator-output kanallarında
zero-shot causal forecast residual üretmek için kullanıldı.

\begin{itemize}
  \item Mechanical recall \(0{,}205\to0{,}390\): Gate B geçti.
  \item Vertical recall \(0{,}096\to0{,}023\): kötüleşti.
  \item Fusion \(0{,}213/23{,}92\) FA/saat: Gate C kaldı.
\end{itemize}

MOMENT foundation model kurulumu Python 3.14 ile eski NumPy/bağımlılık
uyumsuzluğu nedeniyle tamamlanamadı.

\section{Contextual LSTM ve conformal calibration}

\subsection{Amaç}

Önceki modellerin ``büyük residual = anomali'' artefaktını azaltmak için,
doğrudan residual büyüklüğü yerine geçmiş bağlamdan bir sonraki residual
dağılımı tahmin edildi.

\subsection{Girdi ve mimari}

Girdi:

\begin{itemize}
  \item 12 geçmiş adım,
  \item yalnız geçmişten türetilmiş uçuş fazı,
  \item track sin/cos,
  \item gerçek \(\Delta t\) ve cadence,
  \item availability maskeleri,
  \item toplam 19 input boyutu.
\end{itemize}

Tahmin kanalları vertical, speed, heading, east ve north residual'larıdır.
Model bir katmanlı, hidden size 32 LSTM'dir. Her kanal için beklenen ortalama
\(\mu_c\) ve ölçek \(\sigma_c\) üretir; \(\sigma\) yaklaşık
\([0{,}1,5]\) aralığında sınırlandırılır.

Maskeli, kanallar arası açık eşit ağırlıklı Gaussian NLL kullanıldı:

\[
\mathcal{L}
= \frac{1}{\sum_{t,c}m_{t,c}}
\sum_{t,c} m_{t,c}
\left[
\log\sigma_{t,c}
+ \frac{(y_{t,c}-\mu_{t,c})^2}{2\sigma_{t,c}^2}
\right].
\]

Anomaly skoru standartlaştırılmış tahmin hatasıdır:

\[
a_{t,c}
= \frac{|y_{t,c}-\mu_{t,c}|}{\sigma_{t,c}}.
\]

\subsection{Eğitim ve magnitude kontrolü}

\begin{itemize}
  \item 2.929 uçuş,
  \item yaklaşık 1,267 milyon satır,
  \item beş epoch,
  \item NLL \(0{,}795\to0{,}709\).
\end{itemize}

Eğitimli model skorunun magnitude korelasyonu yaklaşık \(0{,}6496\),
karşılaştırma korelasyonu yaklaşık \(0{,}6542\) oldu. Önceki \(0{,}99\)
seviyelerinden daha düşük olduğu için modelin yalnız magnitude kopyalamadığı
kontrolü geçti.

\subsection{Conformal hiyerarşi}

Eşikler channel + phase + cadence düzeyinde kalibre edildi. Grup örneği
1.000'in altındaysa phase, sonra channel düzeyine fallback yapıldı.

Kalibrasyon örneği:

\begin{itemize}
  \item 60/237 part,
  \item 8.182 uçuş,
  \item yaklaşık 3,321 milyon pencere,
  \item yaklaşık 16 milyon uzun-format skor satırı.
\end{itemize}

Alarm bütçesi \(0{,}1,\ 0{,}5,\ 1,\ 2,\ 5\) episode/100 uçuş-saat olarak
ön-kayıtlandı. Bütçe payları yaklaşık S2 \(\%15\), speed \(\%27{,}7\),
heading \(\%22{,}8\), vertical/east/north için \(\%11{,}5\)'er olarak
dağıtıldı. Instant, 30 saniye persistence ve accumulation kararları
değerlendirildi.

\subsection{Contextual truth-v2 sonuçları}

\begin{longtable}{L{0.25\textwidth}L{0.15\textwidth}L{0.18\textwidth}L{0.31\textwidth}}
\toprule
\textbf{Reçete} & \textbf{Bütçe} & \textbf{Recall} & \textbf{Gecikme / not}\\
\midrule
\endfirsthead
\toprule
\textbf{Reçete} & \textbf{Bütçe} & \textbf{Recall} & \textbf{Gecikme / not}\\
\midrule
\endhead
Vertical spike & \(V=5\) & 0,000589 & Yaklaşık 316 s; \(\%0{,}06\) recall.\\
Vertical freeze & \(V=5\) & 0,003061 & Yaklaşık 682 s; \(\%0{,}31\) recall.\\
Speed spike & kayıtlı & 0,056030 & 0 s; en iyi kısa-gecikmeli kanallardan biri.\\
Speed bias & kayıtlı & 0,012826 & Yaklaşık 19,66 s.\\
Heading & kayıtlı & 0,000940 & Yaklaşık 189 s.\\
East/north CUSUM & \(V=0{,}1\) & 0,071268 & Yaklaşık 3.610 s; coverage 0,0286.\\
East/north CUSUM & \(V=5\) & 0,497194 & Yaklaşık 1.101 s; coverage 0,3094.\\
\bottomrule
\end{longtable}

East/north \(V=5\) sonucunda paired clean burden yaklaşık
\(0{,}01998/\text{saat}\) olmasına rağmen aktif enjeksiyonun gözlenebilir
coverage'ı yalnızca \(\%30{,}94\) civarındadır. Altitude dropout, fiziksel
contextual kanal yerine S2 veri kalitesi kapsamındadır.

Debug amacıyla alarm bütçesi ciddi biçimde gevşetilip \(\alpha=0{,}5\)
civarına çıkarıldığında recall bazı reçetelerde \(\%74\)--\(\%92\) seviyesine
ulaşmış; fakat doğal yük yaklaşık 9 alarm/saat olmuştur.

\begin{successbox}[Contextual modelden çıkarılan doğru sonuç]
Model mekanizması tamamen bozuk değildir: gevşek eşikte anomalileri
görebilmektedir ve önceki modeller kadar magnitude bağımlı değildir. Fakat
doğal varyasyon ile anomaliyi kayıtlı operasyonel alarm bütçesinde ayıracak
diskriminasyon oluşmamıştır.
\end{successbox}

Full 237-part kalibrasyon, eşit bütçeli rule karşılaştırması ve ayrı final
blind unseal aşamasına geçilmedi. Kullanılan üç ADS-B günü fit,
calibration/development ve rehearsal rollerinde kullanılmıştır; üç günün
tamamı kör holdout değildir.

\section{Araştırılan fakat tamamlanmayan yöntemler}

\begin{longtable}{L{0.25\textwidth}L{0.61\textwidth}}
\toprule
\textbf{Yöntem} & \textbf{Durum}\\
\midrule
\endfirsthead
\toprule
\textbf{Yöntem} & \textbf{Durum}\\
\midrule
\endhead
Matrix Profile & Literatür ve yol haritasında sağlam unsupervised baseline olarak araştırıldı; resmi tam deney tamamlanmadı.\\
TranAD & Az veri ve adversarial transformer adayı olarak araştırıldı; uygulanmış resmî sonuç yok.\\
Robust covariance / PCA & Baseline/yol haritası seviyesinde kaldı.\\
MOMENT & Python 3.14 ve bağımlılık uyumsuzluğu nedeniyle kurulamadı.\\
RFLY severity/simulation & Preflight veri eksikliği nedeniyle resmi deney tamamlanmadı.\\
ADS-B weighted-loss AE & Önceki adım kapıları geçmediği için koşullu dal çalıştırılmadı.\\
USAD bounded-output smoke & Ana fizik/truth hattı önceliği nedeniyle tamamlanmadı.\\
ADS-B manual top-100 audit & Bekleyen insan incelemesi; ürün kanıtı üretmedi.\\
Full contextual calibration ve blind holdout & Önceki aşamada recall--alarm dengesi geçmediği için açılmadı.\\
\bottomrule
\end{longtable}

\chapter{Kronoloji ve Başarısızlık Zinciri}

\section{Özet kronoloji}

\begin{longtable}{L{0.12\textwidth}L{0.31\textwidth}L{0.47\textwidth}}
\toprule
\textbf{Faz} & \textbf{Hipotez} & \textbf{Sonuç}\\
\midrule
\endfirsthead
\toprule
\textbf{Faz} & \textbf{Hipotez} & \textbf{Sonuç}\\
\midrule
\endhead
ML-0/1 & Tek genel feature uzayı + Isolation Forest & ALFA yaklaşık rastgele, UAV Attack ters skor; monolitik satır modeli reddedildi.\\
ML-1 modüler & Fiziksel modüller feature sulanmasını azaltır & ALFA/UAV iyileşti; küçük validation-normal havuzu eşik kararsızlığı yarattı.\\
ML-2 & LSTM-AE zamansal imzayı öğrenir & UAV Attack'ta IF'yi geçti; ALFA'da önce geçemedi.\\
ML-3 & Ablation, enjeksiyon ve USAD ile hata modu ayrılır & Guidance-only kazancı, platforma özel kalibrasyon ve düşük severity kaçışı ölçüldü; USAD elendi.\\
ML-4 & Daha fazla normal veri LSTM'yi düzeltir & ALFA LSTM-AE \(0{,}731\to0{,}918\); veri hipotezi dar domain için doğrulandı.\\
ML-5/6/7 & Causal özellik ve doğru event onset & Eski CUSUM başarıları ve event recall ciddi düştü; leakage/proxy sorunları düzeltildi.\\
ML-8A & Temporal descriptor + LightGBM daha iyi ayrım sağlar & SEAD ve ALFA'da IF/LSTM baseline'larını geçmedi.\\
ML-9 & EKF ve motor-simetri feature'ları kategori boşluğunu kapatır & Küçük recall artışı; fusion Gate C kaldı.\\
ML-10 & Chronos zero-shot forecast daha güçlü davranış modeli sağlar & Mechanical Gate B geçti; fusion alarm yükü yüzünden kaldı.\\
ML-12 & Tek feature, modül içi seyrelmeyi azaltır & Mechanical recall \(0{,}459\); normal alarm yükü \(38{,}1/\text{saat}\).\\
ML-13 & Sistem ve mekanik kanalları ayırmak fusion kaybını önler & Recall arttı; FA 1,89--3,83 kat yükseldi.\\
ML-14 & Daha fazla ve çeşitli normal session FA'yı azaltır & FA azaldı, recall yaklaşık \(0{,}12\)'ye çöktü.\\
ML-15 & Kalibrasyon drift/jackknife ile genellenir & En iyi \(0{,}1145/8{,}4146\); Gate kaldı.\\
ML-16 & Relative/genlik normalize DL skoru artefaktı çözer & Magnitude korelasyonu düştü; gerçek recall çoğunlukla \(\%6\)'nın altında.\\
RFLY-0 & Yeni gerçek fault verisi Gate C'yi açar & Full-flight proxy geçiyor göründü.\\
RFLY-1 & Gerçek fault interval ve source schema düzeltmesi & Başarı geçersizleşti; corrected FA kabul edilemez.\\
İlk ADS-B & Fiziksel rule + IF sentetik anomalileri düşük yükle yakalar & Truth ve MAD düzeltmesi sonrası \(25{,}537\) yeni alarm/saat; reddedildi.\\
Aktif ADS-B & Robust rule/CUSUM + doğru truth & AUC iyi, doğal alarm yükü çok yüksek.\\
Contextual ADS-B & Bağlam, conformal eşik ve sıkı bütçe ayrımı sağlar & Magnitude sorunu azaldı; kayıtlı bütçede recall sıfıra yaklaştı.\\
\bottomrule
\end{longtable}

\section{Hipotezlerin sıralı olarak reddedilmesi}

\begin{enumerate}
  \item \textbf{Tek genel anomaly modeli}: monolitik IF rastgele/ters sonuç verdi.
  \item \textbf{Fiziksel modüllere bölme}: ALFA'da fayda sağladı, genel transferi çözmedi.
  \item \textbf{Normal veri ekleme}: yanlış alarmı azalttı; sıkı eşikte recall çöktü.
  \item \textbf{Supervised model}: LightGBM, unsupervised baseline'ı geçmedi.
  \item \textbf{Daha karmaşık derin model}: bazı skorlar magnitude/reconstruction artefaktı çıktı.
  \item \textbf{Foundation model}: kategori kazancı operasyonel fusion'a dönüşmedi.
  \item \textbf{CUSUM}: küçük sapmayı topladı, doğal uçuşta alarm saturasyonu oluşturdu.
  \item \textbf{Platform havuzlama}: heterojen normalleri karıştırdı ve genellemeyi bozdu.
  \item \textbf{ADS-B fiziksel tutarlılık}: güçlü sentetik sapmaları yakaladı,
  doğal trafikte alarm yükünü kontrol edemedi.
  \item \textbf{Contextual öğrenme}: magnitude bağımlılığını azalttı; sıkı bütçede
  gerekli recall'ı üretemedi.
\end{enumerate}

\chapter{Bulunan Hatalar ve Uygulanan Düzeltmeler}

\section{Satır, uçuş ve event seviyesinin karıştırılması}

Bir uçuşta çok sayıda satır bulunması bağımsız örnek sayısını artırmaz.
Satır ROC-AUC yüksekken uçuşların çoğu alarm alabilir. Değerlendirme daha
sonra event onset, overlap, uçuş, uçuş-saati, delay ve coverage olarak
ayrıldı.

\section{Tüm uçuşu anomali kabul etmek}

RflyMAD'de klasör etiketi tüm uçuşa uygulanınca model uçuş profili ile fault
aralığını karıştırdı. Gerçek interval kullanıldığında görünen Gate-C başarısı
kayboldu.

\section{Event başlamadan önce açık alarmı tespit saymak}

Eski overlap recall, event öncesi açılmış alarmın event aralığına taşmasını
başarı sayabiliyordu. ``Yeni alarm/onset recall'' tanımı getirildiğinde bazı
sonuçlar yaklaşık \(0{,}59\) overlap'tan \(0{,}19\)--\(0{,}22\) onset
recall'a düştü.

\section{Gelecek bilgisi ve causal CUSUM}

CUSUM center ve \(k\) parametrelerinin tüm seri veya test verisinden
etkilenmesi engellendi. Train-normal baseline sabitlendi ve prefix
invariance testi eklendi. ALFA \code{alt_error_cusum} sonucu
\(0{,}878\to0{,}611\) düştü.

\section{NaN-medyan imputasyonu}

UAV Attack'ta saldırı sırasında kaybolan değer medyana doldurulunca saldırı
daha normal görünüyordu. Missingness ve availability daha sonra özellik ve
maske olarak ele alındı.

\section{Kaynakta olmayan sensörü medyanla üretmek}

Pooled platformlarda bir kaynakta hiç bulunmayan kolonun başka kaynağın
medyanıyla doldurulması engellendi. Source schema maskesi eklendi.

\section{Magnitude artefaktı}

Eğitimli skor, eğitimsiz/shuffled skor ve ham büyüklük arasındaki korelasyon
zorunlu teşhis haline getirildi. Model eğitiminin gerçekten yeni sinyal
öğrenip öğrenmediği bu kontrolle ayrıldı.

\section{Sentetik truth semantiği}

\code{injection_active}, \code{observable_changed} ve
\code{evaluable_truth} ayrıldı. Böylece kodun bir değeri değiştirmeye
çalışması ile model girdisinde fiziksel olarak gözlenebilir değişim oluşması
aynı kabul edilmedi.

\section{Episode merge ile alarm yükünün gizlenmesi}

Uzun süre açık kalan alarmı tek episode saymak alarm sayısını düşük
gösterebilir. Bu nedenle alarm satır coverage'ı ve alarm alan uçuş oranı
episode/saat ile birlikte raporlandı.

\section{Yoğun grid üzerinde hayalî pencere}

Eski temporal descriptor hattında telemetri boşlukları boyunca yoğun 1
saniyelik grid pencereleri üretilebiliyordu. Gap-aware pencereleme ile
düzeltildi; eski çıktı superseded olarak işaretlendi.

\section{Repository paketleme hatası}

Ankorsuz \code{data/} gitignore deseni \code{src/ml/data/} kod paketini de
eşleyebiliyordu. Temiz clone'da import hatası yaratacak bu problem
\code{/data/} biçiminde düzeltilerek ML pipeline modülleri version control
kapsamına alındı.

\chapter{Kök Neden Analizi}

\section{Geçerli ground-truth eksikliği}

ADS-B kaynağında gerçek saldırı veya arıza etiketi yoktur. Sentetik
enjeksiyon, fiziksel anlamı ve feature-level gözlenebilirliği doğrulanmadan
gerçek attack ground-truth'u değildir.

\section{Sensör gözlenebilirliği}

Ping/DoS gibi olaylar mevcut uçuş telemetrisini değiştirmiyorsa telemetri
modeli bu olayı göremez. Ağ paket zamanlaması, sequence number, receiver
durumu veya RF seviyesi olmadan problem gözlenebilir değildir.

\section{Hedef platform uyumsuzluğu}

Aktif ADS-B verisi çoğunlukla insanlı hava araçlarından gelmektedir. UAV
anomali ürününün doğrudan hedef dağılımını temsil etmez.

\section{Domain shift}

ALFA sabit kanat, PX4 multirotor, Rfly ve ADS-B normal davranışları farklıdır.
Aynı residual adı ve eşik bu platformlarda aynı fiziksel anlamı veya
dağılımı taşımamaktadır. Platforma özgü kalibrasyon transfer deneyinde
doğrudan ölçülmüştür.

\section{Bağımsız örnek sayısının küçüklüğü}

Satır hacmi yüksek olsa da bağımsız arıza uçuşu, araç ve session sayısı
düşüktür. Küçük \(n\):

\begin{itemize}
  \item eşik kararsızlığı,
  \item geniş güven aralığı,
  \item fault-family recall'un bir iki örnekle değişmesi,
  \item validation--test normal dağılım kayması
\end{itemize}

oluşturmuştur.

\section{Normal davranışın heterojenliği}

Uçuş fazı, aircraft tipi, manevra, hava, sensör kalitesi, cadence ve provider
semantiği normal dağılımı genişletmektedir. Bu heterojenlik, anomaly ile
normal kuyruğun operasyonel eşikte iç içe geçmesine yol açar.

\section{Recall--yanlış alarm çelişkisi}

Projenin en kararlı sonucu aşağıdaki değiş-tokuştur:

\[
\text{yüksek recall}
\quad\Longrightarrow\quad
\text{yüksek doğal alarm yükü},
\]

\[
\text{düşük doğal alarm yükü}
\quad\Longrightarrow\quad
\text{çok düşük recall}.
\]

ADS-B debug eşiğinde yüksek recall görülebilmesi, mekanizmanın sinyal
üretebildiğini; kayıtlı bütçede recall'ın çökmesi ise doğal varyasyon ile
anomali ayrımının yeterli olmadığını gösterir.

\section{Alarm bütçesi birimi ve ürün sözleşmesi}

Eski deneylerde saat başına 2 veya 12 alarm sınırları kullanılırken aktif
ADS-B çalışması 100 saatte 0,1--5 episode bütçesine geçti. Bütçeyi gevşetmek
``modeli kurtarmak'' değil, farklı alarm yüküne sahip yeni bir ürün
sözleşmesi tanımlamaktır.

\section{Kör final doğrulama eksikliği}

Dar kapsamlı olumlu sonuçlar bulunmasına rağmen tüm genel ürün iddiaları için
ayrı, hiç açılmamış ve yalnız finalde kullanılan bir holdout başarısı
gösterilemedi. Bazı SEAD holdoutları kapalı tutuldu; ADS-B contextual hat ise
blind unseal aşamasına gelmeden önceki kapıda kaldı.

\chapter{Neler Gerçekten Çalıştı?}

\begin{successbox}[Savunulabilir kazanımlar]
\begin{itemize}
  \item Büyük ADS-B verisi streaming, deterministik ve devam ettirilebilir
  biçimde işlendi.
  \item Provenance, checkpoint, checksum, split ve artifact sözleşmeleri
  geliştirildi.
  \item Row, window, event, uçuş ve saat seviyeleri metodolojik olarak ayrıldı.
  \item Leakage, proxy truth, ghost imputation ve magnitude artefaktları
  bulunup ölçüldü.
  \item ALFA'da dar kapsamlı gerçek arıza ayrımı ve LSTM-AE kazanımı gösterildi.
  \item Speed bias ve track freeze gibi güçlü sentetik ADS-B sapmaları
  fiziksel residual ile iyi sıralandı.
  \item S2 katmanı doğal veri kalitesi olaylarını açıklanabilir nedenlerle
  üretebildi.
  \item Contextual LSTM'nin yalnız magnitude kopyalamadığı gösterildi.
  \item Başarısızlık ``model kötü'' ifadesinden çıkarılıp veri, truth,
  gözlenebilirlik ve alarm bütçesi bileşenlerine ayrıldı.
\end{itemize}
\end{successbox}

Bu çıktılar araştırma ve veri mühendisliği başarısıdır. Genel operasyonel
anomali ürününün kabul edildiği anlamına gelmez.

\chapter{Teknik Fizibilite Hükmü}

\section{Hedeflenen ürün}

Hedef; farklı platformlarda, bilinen ve bilinmeyen anomaly türlerinde,
gerçek zamanlı veya yakın gerçek zamanlı çalışabilen, yüksek event recall ve
düşük yanlış alarm üreten genel bir dedektördü.

\section{Neden mevcut koşullarda achievable değil?}

\begin{enumerate}
  \item Hedef ADS-B/UAV saldırıları için yeterli gerçek ve interval etiketli
  veri bulunmamaktadır.
  \item Bazı saldırı etkileri kullanılan sensörlerde gözlenebilir değildir.
  \item ADS-B verisi doğrudan UAV hedef dağılımını temsil etmemektedir.
  \item Platformlar arası normal davranış ve sensör şeması uyumsuzdur.
  \item Bağımsız fault/session örneklemi küçüktür.
  \item Feature ve model iyileştirmeleri doğal normal kuyruğunu anomaly
  kuyruğundan kayıtlı alarm bütçesinde ayıramamıştır.
  \item Olumlu dar sonuçlar genel ve bağımsız final doğrulamaya taşınmamıştır.
\end{enumerate}

\begin{verdictbox}
\textbf{Nihai karar:}

Mevcut problem tanımı değiştirilmeden, yalnız yeni bir model ailesi,
daha büyük sinir ağı, daha uzun eğitim veya hiperparametre taraması eklemek
genel ürün fizibilitesini kanıtlamak için yeterli değildir. Bu nedenle
proje, mevcut veri ve operasyonel kabul şartlarıyla
\textbf{not achievable / teknik olarak fizibl değil} olarak
kapatılmalıdır.
\end{verdictbox}

\section{Bu hükmün sınırı}

Aşağıdaki daha dar ürünler fizibl olabilir:

\begin{itemize}
  \item ADS-B veri kalitesi ve message-gap monitörü;
  \item belirli uçak/receiver grubu için platforma özgü residual monitörü;
  \item yalnız güçlü speed-bias veya track-freeze benzeri sapmaları hedefleyen
  kural dedektörü;
  \item ALFA benzeri sabit platform ve bilinen fault family için özel model;
  \item insan analiste anomaly ranking/triage sağlayan araştırma aracı.
\end{itemize}

Bu ürünler için iddia ve başarı metriği ayrı yazılmalıdır.

\chapter{Yeniden Başlama Koşulları}

Yeni bir çalışmanın aynı başarısızlıkları tekrar etmemesi için aşağıdakiler
başlangıç önkoşulu olmalıdır:

\begin{enumerate}
  \item \textbf{Tek hedef platform}: ADS-B insanlı uçak mı, PX4 UAV mı,
  RF/network attack mı açıkça seçilmeli.
  \item \textbf{Gerçek normal dağılım}: üretim hedefiyle aynı araç, sensör,
  receiver ve cadence'den gelmeli.
  \item \textbf{Event interval etiketi}: anomaly başlangıç ve bitişi
  doğrulanmalı.
  \item \textbf{Gerekli sensör}: DoS için network/RF metadatası, spoofing
  için bağımsız INS/GNSS referansı eklenmeli.
  \item \textbf{Alarm bütçesi}: operasyon ekibiyle önceden dondurulmalı.
  \item \textbf{Blind holdout}: model geliştirmeden önce ayrılmalı ve son
  kabul testine kadar açılmamalı.
  \item \textbf{Platforma özel model}: tek genel ağırlık veya eşik
  varsayılmamalı.
  \item \textbf{Sentetik validasyon}: enjeksiyonun fiziksel anlamı,
  şiddeti, aktif coverage'ı ve gözlenebilirliği doğrulanmalı.
  \item \textbf{Operasyonel metrik}: AUROC/AUPRC yanında event recall,
  onset delay, FA/saat, alarm alan uçuş oranı ve row coverage zorunlu olmalı.
\end{enumerate}

Bu önkoşullar olmadan yeni algoritma denemek, önceki sonucu farklı model
adıyla tekrar etme riski taşır.

\chapter{Beklenen Yönetim Soruları ve Hazır Cevaplar}

\section*{AUC yüksekse neden başarısız sayılıyor?}
\addcontentsline{toc}{section}{AUC yüksekse neden başarısız sayılıyor?}

AUROC yalnız anomaly skorlarının normal skorlardan daha yukarı sıralanma
olasılığını ölçer. Bir eşik seçildiğinde uçuşların \(\%89\)'u alarm alıyorsa
yüksek AUC operasyonel başarı değildir.

\section*{Neden Isolation Forest kullanıldı?}
\addcontentsline{toc}{section}{Neden Isolation Forest kullanıldı?}

Gerçek saldırı etiketi olmayan normal-only novelty problemine uygun, hızlı
ve açıklanabilir bir baseline olduğu için kullanıldı. Ancak birçok deneyde
ham residual magnitude dedektörüne dönüştüğü ölçüldü.

\section*{Neden LightGBM kullanıldı?}
\addcontentsline{toc}{section}{Neden LightGBM kullanıldı?}

Etiketli event pencerelerindeki doğrusal olmayan descriptor ilişkilerini
öğrenmek ve unsupervised modelin kaçırdığı fault ailelerini yakalamak için.
SEAD ve ALFA'da kayıtlı baseline'ı geçmedi.

\section*{Neden derin öğrenme kullanıldı?}
\addcontentsline{toc}{section}{Neden derin öğrenme kullanıldı?}

Zamansal bağımlılığı, uçuş dinamiğini ve doğrusal olmayan normal davranışı
öğrenmek için. ALFA'da normal veri artırılınca fayda sağladı; heterojen
SEAD/ADS-B ortamında magnitude ve reconstruction artefaktı baskın kaldı.

\section*{Daha büyük model sorunu çözmez mi?}
\addcontentsline{toc}{section}{Daha büyük model sorunu çözmez mi?}

Sensörde bulunmayan sinyali, hatalı etiketi veya doğal/anomaly dağılımlarının
operasyonel eşikte çakışmasını model boyutu çözmez. Daha büyük model
overfitting ve kalibrasyon riskini artırabilir.

\section*{Daha fazla veri çözmez mi?}
\addcontentsline{toc}{section}{Daha fazla veri çözmez mi?}

Aynı tür etiketsiz heterojen verinin artması doğal alarm kuyruğunu daha iyi
ölçer; fakat gerçek anomaly etiketi ve eksik sensörü üretmez. ALFA örneği,
doğru domain normal verisinin dar bir problemde faydalı olduğunu; SEAD
örneği ise heterojen normal artışının recall'ı kurtarmadığını göstermiştir.

\section*{Eşiği gevşetsek olmaz mı?}
\addcontentsline{toc}{section}{Eşiği gevşetsek olmaz mı?}

Recall yükselir; fakat deneylerde yaklaşık 9--25 alarm/saat veya daha yüksek
yük oluşmuştur. Bu, mevcut ürün şartını karşılamak değil alarm sözleşmesini
değiştirmektir.

\section*{ALFA sonucu başarı değil mi?}
\addcontentsline{toc}{section}{ALFA sonucu başarı değil mi?}

Dar ve platforma özel bir araştırma başarısıdır. Küçük örneklemli, bazı
sınıflarda \(n=1\)--4 düzeyinde, final blind holdout'suz ve diğer platformlara
transfer edilemeyen sonuç genel ürün kanıtı değildir.

\section*{Sentetik anomaly neden yeterli değil?}
\addcontentsline{toc}{section}{Sentetik anomaly neden yeterli değil?}

Enjeksiyon gerçek fiziksel saldırının etkisini, şiddetini ve sensör
gözlenebilirliğini temsil etmelidir. Aksi halde model gerçek saldırıyı değil,
bizim yazdığımız dönüşüm kodunu öğrenir.

\section*{ADS-B ile drone anomalisi tespit edildi mi?}
\addcontentsline{toc}{section}{ADS-B ile drone anomalisi tespit edildi mi?}

Hayır. Aktif veri büyük ölçüde insanlı hava aracı ADS-B trafiğidir ve
doğrulanmış UAV ground-truth'u yoktur. Yapılan çalışma ADS-B fiziksel
tutarlılık ve veri kalitesi analizidir.

\section*{Project tamamen imkânsız mı?}
\addcontentsline{toc}{section}{Project tamamen imkânsız mı?}

Evrensel olarak değil. Ancak mevcut kapsam, veri, gözlem kanalları ve alarm
bütçesi altında hedeflenen genel operasyonel ürün achievable değildir. Dar
ve yeniden tanımlanmış alt ürünler mümkündür.

\section*{Neden daha fazla yöntem denenmedi?}
\addcontentsline{toc}{section}{Neden daha fazla yöntem denenmedi?}

Matrix Profile, TranAD, PCA, robust covariance ve MOMENT gibi adaylar
araştırıldı. Fakat temel truth, gözlenebilirlik ve alarm yükü sorunları model
ailesinden bağımsızdır. Önceki kapılar geçilmeden sınırsız yöntem taramak
sonuç sonrası seçim yanlılığı yaratır.

\chapter{Sonuç}

Bu proje yalnızca ``denedik, olmadı'' seviyesinde kapanmamıştır. Farklı veri
setleri, fiziksel varsayımlar, ML ve DL modelleri, kalibrasyon yöntemleri,
sentetik truth sözleşmeleri ve operasyonel alarm kararları ardışık biçimde
test edilmiştir.

İlk bakışta olumlu görünen sonuçlar leakage, proxy etiket, tüm-uçuş etiketi,
magnitude artefaktı veya alarm birimi açısından tekrar denetlenmiştir.
Düzeltmelerden sonra genel örüntü değişmemiştir:

\begin{verdictbox}
Kabul edilebilir yanlış alarm seviyesinde recall operasyonel olarak yetersiz;
recall yükseltildiğinde ise alarm yükü kabul edilemezdir. Mevcut veride gerçek
anomaly ground-truth'u, hedef platform temsili ve gerekli sensör
gözlenebilirliği de eksiktir. Bu nedenle genel operasyonel anomali tespit
ürünü mevcut şartlarla teknik olarak fizibl değildir.
\end{verdictbox}

\appendix

\chapter{Repo İçindeki Ana Kanıt Belgeleri}

\begin{longtable}{L{0.34\textwidth}L{0.56\textwidth}}
\toprule
\textbf{Belge} & \textbf{İçerik}\\
\midrule
\endfirsthead
\toprule
\textbf{Belge} & \textbf{İçerik}\\
\midrule
\endhead
\path{docs/decisions.md} & Aktif ve tarihsel mimari/deney kararları, ADR kayıtları.\\
\path{docs/adsb_overall_model_report_2026-07-14.md} & Aktif ADS-B fizik, kural, DL ve doğal yük sonuçları.\\
\path{docs/DURUM_TESHIS_VE_YOL_HARITASI.md} & Durum teşhisi ve aşamalı yol haritası.\\
\path{docs/adsb_contextual_candidate_v1_prereg_2026-07-14.md} & Contextual LSTM ön-kayıt sözleşmesi.\\
\path{docs/adsb_contextual_physics_v1_alarm_budget_prereg_2026-07-14.md} & ADS-B alarm bütçesi ve Pareto sözleşmesi.\\
\path{archive/2026-07-10_legacy_non_adsb_ml/docs/ML_YETERSIZLIKLER_KAYDI.md} & Eski ML hattının yetersizlikleri ve açık sınırlar.\\
\path{archive/2026-07-10_legacy_non_adsb_ml/docs/ML1_BULGULAR_VE_HATALAR.md} & ML fazlarının bulgu ve hata kayıtları.\\
\path{archive/2026-07-10_legacy_non_adsb_ml/docs/ANOMALI_METOT_ARASTIRMASI.md} & Yöntem araştırması ve hata modu--metot eşlemesi.\\
\path{archive/2026-07-10_rejected_adsb_attempts/docs/CODEX_ADSB_BEHAVIORAL_V1_RESULTS.md} & Reddedilen ilk ADS-B behavioral yaklaşım sonuçları.\\
\bottomrule
\end{longtable}

\chapter{Kısa Terimler Sözlüğü}

\begin{longtable}{L{0.21\textwidth}L{0.69\textwidth}}
\toprule
\textbf{Terim} & \textbf{Açıklama}\\
\midrule
\endfirsthead
\toprule
\textbf{Terim} & \textbf{Açıklama}\\
\midrule
\endhead
ADS-B & Uçağın kimlik, konum, hız ve durum bilgisini yayınladığı gözetim sistemi.\\
Residual & Ölçülen/bildirilen değer ile fiziksel veya öğrenilmiş beklenen değer arasındaki fark.\\
AUROC & Pozitif örneklerin negatiflerden daha yüksek sıralanma olasılığına dayalı ranking metriği.\\
AUPRC & Precision--recall eğrisi altındaki alan; sınıf dengesizliğine AUROC'dan daha duyarlı.\\
Event recall & Gerçek anomaly olaylarının ne kadarında geçerli alarm oluştuğu.\\
False alarm/hour & Normal sürede oluşan yanlış alarm sayısının saatle normalize edilmiş hali.\\
CUSUM & Küçük fakat kalıcı sapmaları zaman içinde biriktiren ardışık değişim skoru.\\
Isolation Forest & Normal dağılımın dışında kalan örnekleri ağaç izolasyonu ile skorlayan yöntem.\\
Autoencoder & Normal veriyi yeniden üretmeyi öğrenen; reconstruction error'ı anomaly skoru yapan sinir ağı.\\
Conformal calibration & Kalibrasyon skorlarının ampirik dağılımından kapsama kontrollü eşik üreten yaklaşım.\\
Magnitude artefaktı & Model skorunun öğrenilmiş davranış yerine büyük ham sinyal genliğini takip etmesi.\\
Proxy truth & Gerçek anomaly interval'ı yerine yaklaşık klasör, uçuş veya enjeksiyon durumu etiketi kullanılması.\\
Observable coverage & Enjeksiyon aktifken kullanılan model girdilerinde gerçekten gözlenebilir değişim oluşan oran.\\
Blind holdout & Model ve eşik seçimi tamamlanana kadar hiç açılmayan final değerlendirme bölümü.\\
\bottomrule
\end{longtable}

\chapter{Overleaf Derleme Notu}

Bu dosya tek başına derlenebilir. Overleaf üzerinde:

\begin{enumerate}
  \item \code{ml_anomali_tespiti_fizibilite_final_raporu_2026-07-16.tex}
  dosyasını yeni projeye yükleyin.
  \item \textbf{Menu} bölümünde \textbf{Compiler} seçeneğini
  \textbf{XeLaTeX} yapın.
  \item \textbf{Recompile} düğmesine basın.
\end{enumerate}

Rapor dış görsel veya ayrı bibliography dosyasına ihtiyaç duymaz.

\end{document}
